\chapter{Introduction}

While players primarily experience gameplay and immersive worlds, these experiences are often driven - behind the scenes - by complex algorithms and mathematical concepts.

As the creative and technical vision of the project expand, the developers' reliance on advanced mathematics scales with it; whether it be to achieve more realistic physics,
procedural generation, or optimize existing systems.
However, when developers turn to contemporary research literature to implement or optimize these visions, they may face a critical roadblock.
A significant gap exists between the formal academic notation that can be found in scientific papers, such as summation ($\sum$), or product ($\prod$) symbols,
and the practical, programming-orientated skill-set of developers who may lack an existing formal background in these fields of academia.
Consequently, the literature that has been created to provide a way to establish and solve problems become arguably inaccessible to a large portion of the development community.

This report aims to measure a wide-range of developers' baseline compression of different common-place concepts and notation that can be found in mathematic literature, and
comparing it to the level of explanation and context that can be found in these papers.

\section{Defining the developer}

In the context of this report, a `developer' is someone who:

\begin{enumerate}
    \item has written code in \CPP \textit{or in a comparable game-engine language such as} \CS
    \item has less than \textbf{2 years} of formal further/higher-education in mathematics
    \item is currently engaged in either \textbf{independent, hobbyist, or industry} game development \textit{or currently engaged in similar fields such as software engineering or web development}
\end{enumerate}

\pagebreak

\section{Defining the paper}

In the context of this report, a `paper' may refer to:

\begin{enumerate}
    \item A report that is written specifically with game development in mind
    \item A report that is written with physicists or mathematicians in mind
    \item A report from adjacent technical fields, such as computer science, graphics programming or artificial intelligence, 
            which uses any mathematical concepts.
    \item Any peer-reviewed academic literature, journal article, or any other formal academic writing that contains contemporary physics or mathematical content.
\end{enumerate}